% Programming with InitechBase, Chapter 21: Searching (AN-16).  Panel
% replayed (examples/p21-countfor.dot); sources cmd/query.c (LOCATE,
% CONTINUE, FOUND), the program COUNTFOR.PRG.
\chapter{Searching}
\Gates{test-pubs-examples, test-samir-query}
\cmd{SEEK} finds a key in an index at once.  A search on a condition that no
index serves -- employees hired before a date, say -- is done with
\cmd{LOCATE} and \cmd{CONTINUE}.

\begin{ChapterCovers}
\item A search loop on \fn{FOUND()}
\item Counting the records found
\end{ChapterCovers}

\ATEHead{Preparing for This Chapter}
Start InitechBase with a fresh copy of the sample diskette, and type
\type{USE STAFF}.

\ATEHead{A Search Loop}
\cmd{LOCATE} finds the first record that meets its condition and makes
\fn{FOUND()} true; each \cmd{CONTINUE} finds the next.  When there is none,
\fn{FOUND()} is false.  A loop on \fn{FOUND()} therefore visits every record that
meets the condition:
\ProgramListing{COUNTFOR.PRG}
\begin{Steps}
\item Type \type{MBEFORE = CTOD('01/01/83')}.
\item Type \type{DO COUNTFOR}.
\end{Steps}
\DotExample{examples/p21-countfor}
The condition of the \cmd{LOCATE} uses the memory variable \texttt{MBEFORE}, so the
same program searches for any date.

\ATEHead{Where Do You Go From Here?}
Chapter~22 puts the Initech Accounting System together.

\ATEHead{Summary}
\begin{enumerate}
\item \cmd{LOCATE FOR} \xp{condition}, then \cmd{CONTINUE} in a loop on
  \fn{FOUND()}, visits every record that meets the condition.
\item The condition may use memory variables.
\end{enumerate}
